\chapter*{Resumen}
\addcontentsline{toc}{chapter}{Resumen}

Este Trabajo Final de Ingeniería en Informática aborda el desarrollo y validación de un asistente virtual inteligente diseñado específicamente para la preparación del examen de certificación Project Management Professional (PMP) del Project Management Institute (PMI). Ante la complejidad del examen y las limitaciones de los métodos de estudio tradicionales —a menudo estáticos, costosos o carentes de personalización—, se propone una solución tecnológica basada en Modelos de Lenguaje Grande (LLM) integrados con técnicas de Generación Aumentada por Recuperación (RAG).

La metodología empleada combina el rigor de la ingeniería de software moderna con principios pedagógicos adaptativos. El sistema, construido sobre una arquitectura web escalable utilizando Next.js y Python/FastAPI, implementa un motor de razonamiento capaz de explicar conceptos abstractos, evaluar el conocimiento del usuario mediante simulación de escenarios situacionales y ofrecer retroalimentación socrática inmediata.

Los resultados obtenidos validan la hipótesis de que la inteligencia artificial generativa, cuando se diseña con controles éticos y precisión técnica, puede actuar eficazmente como un tutor personalizado. Las pruebas con usuarios demostraron una mejora significativa en la comprensión de conceptos clave de gestión de proyectos y una reducción en la ansiedad asociada a la preparación del examen. Este trabajo no solo entrega una herramienta funcional, sino que establece un precedente sobre cómo la tecnología puede democratizar el acceso a certificaciones profesionales de alto nivel.

\vspace{1cm}
\textbf{Palabras clave:} Project Management Professional, Inteligencia Artificial Generativa, Tutor Inteligente, Educación Profesional, Ingeniería de Software.